\documentclass[12pt, pdf, xcolor={table, dvipsnames}, paperheight=8cm,paperwidth=12cm]{beamer}
\definecolor{tugreen}{RGB}{132,184,25}
\usepackage{assets/themes/beamerthemeTUDortmund}
\usepackage{assets/themes/beamercolorthemetudortmund}

%\expandafter\def\csname ver@transparent.sty\endcsname{}
%\usepackage[edges]{forest}
%\usepackage{pgf-umlsd}
%\usepackage{dsfont}
\usepackage[utf8]{inputenc}
\usepackage[ngerman]{babel}
\usepackage[style=german,babel]{csquotes}
\usepackage{glossaries}
\makeindex
\usepackage{tabularx}
\usepackage{multirow}
%\usepackage{amsmath}
%\usepackage{graphicx}
%\usepackage{pgfplots} 
%\usepackage{fontawesome}
%\usepackage{import}
%\usepackage{listings}
%\usepackage{titleref}
%\usepackage{caption}
%\usepackage{tikz}
%\usepackage{../assets/tikz-uml}
%\usepackage{biblatex}
%\addbibresource{literatur.bib}
%\usetikzlibrary{backgrounds,shapes.multipart}
%\pgfdeclarelayer{myback}
%\pgfsetlayers{background,main,myback}
%\makeatletter
%\newcommand*{\currentname}{\TR@currentTitle}
%\newcommand*{\overlaynumber}{\number\beamer@slideinframe}
%\makeatother
%\usetikzlibrary{positioning,arrows,automata,shapes.multipart,calc, fit, decorations.shapes,shadows}
%\tikzset{
%  mylabel/.style = {font=\footnotesize, midway, fill=white, anchor=center},
%  decorate sep/.style 2 args=
%    {decorate,decoration={shape backgrounds,shape=circle,shape size=#1,shape sep=#2}},
%  invisible/.style={opacity=0, text opacity=0},
%  visible on/.style={alt={#1{}{invisible}}},
%  alt/.code args={<#1>#2#3}{%
%    \alt<#1>{\pgfkeysalso{#2}}{\pgfkeysalso{#3}} % \pgfkeysalso doesn't change the path
%  },
%  onslide/.code args={<#1>#2}{%
%  \only<#1>{\pgfkeysalso{#2}} % \pgfkeysalso doesn't change the path
%    }
%}
%\forestset{
%  step/.style={
%    /tikz/visible on={<+->},
%    edge={/tikz/visible on={<+->}},
%    for descendants={
%        /tikz/visible on={<.->},
%        edge={/tikz/visible on={<.->}}
%    }
%  },
%  highlight/.style={
%    for tree={
%      /tikz/onslide=<.>{draw=red},
%      edge+={/tikz/onslide=<.>{draw=red}}
%    }
%  }
%}
%
\author{Dennis Ochocki}
\title[]{Sensordatenverarbeitung durch metamodelierte best-of-breed Lösung}
%\subtitle{Grafana, Prometheus und Docker}
%\titlegraphic{
%\includegraphics[width=0.35\textwidth]{../assets/themes/img/tud_logo_rgb}\hspace*{3.75cm}~%
%   \includegraphics[width=0.25\textwidth]{../assets/themes/img/fakultaet.pdf}
%}
%\logo{}
%\institute{}
%\date{\today}
%\subject{}
%\setbeamertemplate{navigation symbols}{}
%\setbeamertemplate{frametitle}{%
%    \nointerlineskip%
%    \begin{beamercolorbox}[wd=\paperwidth,ht=2.6ex,dp=1.6ex]{frametitle}
%        \hspace*{1ex}\insertframetitle%
%    \end{beamercolorbox}%
%}
%\setlength{\skip\footins}{0cm}
%\addtobeamertemplate{frametitle}{}{%
%\begin{tikzpicture}[remember picture,overlay]
%\node[anchor=north east,yshift=0pt] at (current page.north east) {\includesvg[height=0.7cm]{../themes/img/fakultaet}};
%\end{tikzpicture}}
%
\AtBeginSection[]{
  \begin{frame}
  \vfill
  \centering
  \begin{beamercolorbox}[sep=8pt,center,shadow=false,rounded=false]{title}
    \usebeamerfont{title}\insertsectionhead\par%
  \end{beamercolorbox}
  \vfill
  \end{frame}
}


\newglossaryentry{latex}
{
        name=latex,
        description={Is a mark up language specially suited for 
scientific documents}
}

\newglossaryentry{maths}
{
        name=mathematics,
        description={Mathematics is what mathematicians do}
}

\newglossaryentry{formula}
{
        name=formula,
        description={A mathematical expression}
}

\newacronym{iot}{IoT}{Internet of things}
\newacronym{pql}{PQL}{Prometheus Query Language}
\newacronym{tsdb}{TSDB}{Time Series Database}
\newacronym{promql}{PromQL}{Prometheus Query Language}
\newacronym{bzgl}{bzgl.}{bezüglich}
\newacronym{IP}{IP}{Internet Protocol}
\newacronym{IPv4}{IPv4}{Internet Protocol version 4}
\newacronym{IPv6}{IPv6}{Internet Protocol version 6}
\newacronym{ip-address}{IP-address}{Internet Protocol Address}
\newacronym{http}{HTTP}{Hyper Text Transfer Protocol}
\newacronym{ebnf}{EBNF}{Extended Backus–Naur form}
\newacronym{json}{JSON}{JavaScript Object Notation}
\newacronym{tudortmund}{TU Dortmund}{Technische Universität Dortmund}
\newacronym{ide}{IDE}{integrated development environment}
\newacronym{dsl}{DSL}{domain specific language}
\newacronym{mgl}{MGL}{meta graph language}
\newacronym{msl}{MSL}{meta style language}
\newacronym{w}{W}{Watt}
\newacronym{percent}{\%}{percent}
\newacronym{degreeC}{$^\circ$C}{degree celcius}
\newacronym{url}{URL}{Uniform Resource Locator}
\newacronym{ui}{UI}{User Interface}
\newacronym{emf}{EMF}{Eclipse Modeling Framework}
\newacronym{uml}{UML}{Unified Modeling Language}

\begin{document}
	
\maketitle

\begin{frame}{Übersicht}
	\tableofcontents
\end{frame}

\section{Motivation}
\begin{frame}{Motivation}
	Smart Home ist immer präsenter und dank \gls{iot} erst möglich. Immer mehr Geräte, sog. \enquote{Smart Devices}, kommen auf den Markt, die in solche Systeme integriert werden können. 
	
	
	\begin{block}{Was ist Smart Home:}
		\begin{itemize}
			\item Viele verschiedene Geräte in einem Netzwerk
			\item Jedes Gerät ist Beobachter und/oder Akteur
			\item Dient der Qualitätssteigerung der vernetzten Örtlichkeit
		\end{itemize}
	\end{block}
\end{frame}

\begin{frame}{Motivation}
Meist sind die Anwendungsfälle trivial und einfach zu realisieren, damit schöpft man aber nicht alles an Funktionalität aus, was möglich wäre.
\begin{block}{Was sind die Schwierigkeiten?}
	\begin{itemize}
		\item Bei vielen Smart Devices ist es schwer den Überblick zu behalten
		\item Assoziationen von Geräte-IP zur physikalischer Position muss gemerkt werden
		\item Muss nicht auf eine Wohnung begrenz sein, mehrere Örtlichkeiten, Häuser Städte möglich.
	\end{itemize}
\end{block}
\end{frame}


\section{Grundlagen}

\begin{frame}{Senoren und Akteure}
	Sensoren und Akteure sind die atomaren Bestandteile eines solchen Netzwerkes.
	
	\vspace{1em}
	Sensoren liefern Daten/Metriken und Aktuere können auf das Umfeld einwirken um damit auf die erfassten Werte zu reagieren.
	\begin{block}{Beispiel: Sensor}
		\vspace{0.5em}
		\begin{tabularx}{\linewidth}{lX}Thermometer: & Erfasst Temperatur, stellt Werte zur  Verfügung
		\end{tabularx}
	\end{block}
	\begin{block}{Beispiel: Akteur}
	\vspace{0.5em}
	\begin{tabularx}{\linewidth}{lX}Thermostat: &Steuert Heizung, hat Einfluss auf die Temperatur
	\end{tabularx}
	\end{block}
Ein Gerät kann ein Sensor und/oder ein Akteur sein.
\end{frame}

\begin{frame}{Metrik}
	Eine Metrik ist ein 
\end{frame}

\begin{frame}{Prometheus}
	\dots ist eine Systemüberwachungs- und Benachrichtigungssoftware
	
	Es greift nach Prometheusdefinition standartisierte Schnittstellen der Sensoren ab um darüber Messwerte/Metriken zu lesen und in der internen Time Series DB abzuspeichern.
	
	Über die \gls{pql} kann man auf diese Daten zugreifen.
\end{frame}

\begin{frame}{Prometheus - Sensor}
	Sensoren können hier angebunden werden und liefern nach besagtem Standart folgende Ausgabe:
	
	
\end{frame}

\begin{frame}{Grafana}
	\dots dient dem Darstellen und der Analyse der durch z.B. Prometheus gewonnenen Daten.
	
	Alternativ zu Prometheus können auch andere Datenquellen wie z.B. JSON, Cassandra, MySQL, ... angebunden werden.
	
	Es können verschiedene Darstellungen gewählt werden (Tabelle, Graph, Säulendiagramm)
\end{frame}

\begin{frame}{Cinco}
\end{frame}

\section{Technologien}
\begin{frame}{CNNM}
\end{frame}


% \section{Motivation}
% \begin{frame}{\secname}
% \begin{exampleblock}{Konzept}
% \onslide<+->
% \begin{itemize}[opacity=.2]
% \item<+-> Haben reelle Situation aus dem Straßenverkehr
% \item<+-> Wollen diese so einfach wie möglich textuel beschreiben
% \item<+-> Gängige Formate wie z.B. OpenScenario sehr komplex
% \item<+-> Hoher Grad von Abstraktion gewünscht
% \item<+-> Geht das überhaupt ohne essentiellen Informationsverlust?
% \end{itemize}
% \end{exampleblock}
% \end{frame}

\begin{frame}[plain,allowframebreaks]{Verwandte Arbeiten}
	%\nocite{*}
	%\printbibliography[heading=none, keyword={related}]
\end{frame}

\begin{frame}[plain,allowframebreaks]{Mögliche Literatur}
	%\printbibliography[heading=none, keyword={possible}]
\end{frame}
\end{document}